\section{CPU 架构设计说明书}

\subsection{本章对应的赛题要求}

本章直接对应赛题中的“CPU 架构详细设计”和“CPU 架构设计说明书”，重点覆盖以下内容：
\begin{itemize}
    \item 五级流水结构、流水寄存器接口与控制信号组织；
    \item IF/ID/EX/MEM/WB 五级功能划分，以及寄存器堆、ALU、访存、写回等核心模块职责；
    \item 数据冒险、控制冒险、结构冒险的当前实现路径与工程边界；
    \item 参数化与扩展性设计，含 \texttt{XLEN}、同步存储路径和最小 SoC 承接方式。
\end{itemize}

\subsection{系统总体架构}

\texttt{YH\_rv\_cpu} 采用“CPU 核 + SoC 封装 + FPGA 原型顶层”的层次化组织方式。CPU 内核负责指令执行、
控制流处理和异常响应；SoC 层负责连接同步指令存储、数据存储和最小 MMIO 外设；
FPGA 顶层负责将工程映射到 Nexys A7-100T 的原型实现路径。
\begin{figure}[H]
    \centering
    \begin{tikzpicture}[node distance=8mm and 10mm, every node/.style={font=\small}]
        \tikzset{
            archblock/.style={
                draw=black!70,
                rounded corners=2pt,
                fill=white,
                align=center,
                inner sep=4pt
            },
            archcore/.style={
                archblock,
                fill=blue!5,
                minimum width=43mm,
                minimum height=38mm
            },
            archif/.style={
                archblock,
                fill=green!5,
                minimum width=39mm,
                minimum height=15mm
            },
            archfabric/.style={
                archblock,
                fill=gray!10,
                minimum width=24mm,
                minimum height=36mm
            },
            archmem/.style={
                archblock,
                fill=green!5,
                minimum width=39mm,
                minimum height=15mm
            },
            archmmio/.style={
                archblock,
                fill=orange!8,
                minimum width=39mm,
                minimum height=15mm
            },
            archlink/.style={
                <->,
                >=Latex,
                semithick,
                draw=black!75
            },
            archgroup/.style={
                draw=black!40,
                rounded corners=3pt,
                inner sep=8pt
            }
        }

        \node[archcore] (cpu) {
            \textbf{CPU Core}\\
            \scriptsize \texttt{YH\_rv\_cpu}\\
            \scriptsize IF / ID / EX / MEM / WB\\
            \scriptsize Hazard / Forwarding\\
            \scriptsize CSR / Trap / Redirect
        };

        \node[archif, right=28mm of cpu, yshift=14mm] (imem) {
            \textbf{Instruction ROM}\\
            \scriptsize \texttt{YH\_rv\_sync\_imem\_rom}\\
            \scriptsize Sync ifetch path
        };

        \node[archfabric, right=10mm of cpu, yshift=-12mm] (fabric) {
            \textbf{Address Decode}\\
            \scriptsize load/store dispatch\\
            \scriptsize MMIO routing
        };

        \node[archmem, right=14mm of fabric, yshift=8mm] (dmem) {
            \textbf{Data RAM}\\
            \scriptsize \texttt{YH\_rv\_dmem\_ram}\\
            \scriptsize Load / Store path
        };

        \node[archmmio, right=14mm of fabric, yshift=-11mm] (mmio) {
            \textbf{MMIO Peripherals}\\
            \scriptsize UART / TIMER / DONE
        };

        \draw[archlink] ([yshift=10mm]cpu.east) -- node[above, font=\scriptsize]{ifetch interface} (imem.west);
        \draw[archlink] ([yshift=-9mm]cpu.east) -- node[above, font=\scriptsize]{load / store / mmio} (fabric.west);
        \draw[archlink] (fabric.east |- dmem.west) -- node[above, font=\scriptsize]{data memory} (dmem.west);
        \draw[archlink] (fabric.east |- mmio.west) -- node[below, font=\scriptsize]{peripheral access} (mmio.west);

        \node[archgroup, fit=(cpu)(imem)(fabric)(dmem)(mmio),
              label={[font=\small]above:SoC 封装（\texttt{YH\_rv\_cpu\_soc}）}] (soc) {};
        \node[archgroup, thick, inner sep=16pt, fit=(soc),
              label={[font=\small]above:FPGA 原型顶层（\texttt{YH\_rv\_cpu\_fpga\_top}）}] (fpga) {};
    \end{tikzpicture}
    \caption{系统总体模块框图}
\end{figure}


图~1 强调了三层边界关系：CPU 内核承接执行语义，SoC 层负责可运行的软件系统环境，
FPGA Top 负责将同一套工程映射到原型平台。这种组织方式避免了“单独为说明书维护一套结构、
单独为原型实现再维护另一套结构”的割裂问题，有利于保证设计说明、RTL、验证和实现报告之间的一致性。

\begin{table}[H]
    \centering
    \small
    \caption{核心模块职责}
    \begin{tabularx}{0.95\textwidth}{L{0.28\textwidth} Y}
        \toprule
        模块 & 作用 \\
        \midrule
        \texttt{YH\_rv\_cpu} & CPU 顶层，组织 IF/ID/EX/MEM/WB 五级流水、前递、redirect、CSR 与 trap \\
        \texttt{YH\_rv\_cpu\_hazard\_unit} & 负责 load-use hazard 检测和前递选择，是停顿控制的关键模块 \\
        \texttt{YH\_rv\_cpu\_soc} & SoC 封装层，连接同步指令存储、数据存储与最小 MMIO 外设 \\
        \texttt{YH\_rv\_sync\_imem\_rom} & 服务同步取指路径，使 CPU 与 FPGA 原型实现保持一致 \\
        \texttt{YH\_rv\_dmem\_ram} & 负责 load/store 数据访问，并与 MMIO 共同组成访问路径 \\
        \texttt{YH\_rv\_cpu\_fpga\_top} & FPGA 原型顶层与综合参数入口，用于承接 \texttt{impl50} 等实现流程 \\
        \bottomrule
    \end{tabularx}
\end{table}

\subsection{CPU 五级流水组织}

处理器采用经典五级流水：
\begin{enumerate}
    \item IF：同步取指、PC 更新与 redirect 接收；
    \item ID：译码、寄存器读、立即数生成和控制信号准备；
    \item EX：ALU 运算、分支跳转判断、CSR 访问以及 trap/\texttt{mret} 决策；
    \item MEM：load/store 与 MMIO 访问；
    \item WB：将 ALU、访存或 CSR 结果回写寄存器堆。
\end{enumerate}

这种结构的优势在于边界清楚、验证成本可控，并且足以支撑 CoreMark、
\texttt{riscv-tests} 和最小 SoC 软件运行。初赛版本没有引入 cache、多发射或乱序机制，
而是优先保证结构可说明、行为可验证、结果可复现。
\begin{figure}[H]
    \centering
    \begin{tikzpicture}[node distance=6mm and 6mm, every node/.style={font=\small}]
        \node[reportbox, minimum width=24mm, minimum height=13mm, fill=black!2] (if) {\textbf{IF}\\\scriptsize PC / 取指};
        \node[reportbox, minimum width=24mm, minimum height=13mm, fill=black!2, right=of if] (id) {\textbf{ID}\\\scriptsize 译码 / 寄存器};
        \node[reportbox, minimum width=24mm, minimum height=13mm, fill=black!2, right=of id] (ex) {\textbf{EX}\\\scriptsize ALU / 分支 / CSR};
        \node[reportbox, minimum width=24mm, minimum height=13mm, fill=black!2, right=of ex] (mem) {\textbf{MEM}\\\scriptsize Load / Store};
        \node[reportbox, minimum width=24mm, minimum height=13mm, fill=black!2, right=of mem] (wb) {\textbf{WB}\\\scriptsize 写回};

        \draw[reportline] (if) -- (id);
        \draw[reportline] (id) -- (ex);
        \draw[reportline] (ex) -- (mem);
        \draw[reportline] (mem) -- (wb);

        \node[below=1.5mm of $(if.east)!0.5!(id.west)$, font=\scriptsize, text=black!55] {IF/ID};
        \node[below=1.5mm of $(id.east)!0.5!(ex.west)$, font=\scriptsize, text=black!55] {ID/EX};
        \node[below=1.5mm of $(ex.east)!0.5!(mem.west)$, font=\scriptsize, text=black!55] {EX/MEM};
        \node[below=1.5mm of $(mem.east)!0.5!(wb.west)$, font=\scriptsize, text=black!55] {MEM/WB};

        \node[reportbox, fill=green!8, minimum width=58mm, minimum height=13mm, below=18mm of id] (haz) {
            \textbf{相关处理控制单元}\\
            \scriptsize load-use 检测、\texttt{stall\_decode} 与前递选择
        };

        \draw[reportdash] (haz.north) -- node[left, font=\scriptsize]{stall\_decode} (id.south);
        \draw[reportdash] (haz.north west) to[out=125,in=-120] node[left, font=\scriptsize]{停顿控制} ([xshift=-1mm]if.south);
        \draw[reportdash] (haz.north east) to[out=55,in=-110] node[right, font=\scriptsize]{前递选择} ([xshift=2mm]ex.south east);
        \draw[reportdash, blue!70!black] (mem.south) to[out=-110,in=-15] node[below, font=\scriptsize]{EX/MEM 前递} ([xshift=1mm]ex.south east);
        \draw[reportdash, blue!70!black] (wb.south) to[out=-110,in=-5] node[below, font=\scriptsize]{MEM/WB 前递} ([xshift=5mm]ex.south east);
        \draw[reportline, red!75!black] ([yshift=1mm]ex.north) to[out=100,in=80] node[above, font=\scriptsize]{redirect / trap / mret} ([yshift=1mm]if.north);
    \end{tikzpicture}
    \caption{五级流水与关键控制通路示意图}
\end{figure}


图~2 突出展示了两类关键通路：一类是 EX 级返回 IF 的控制流重定向路径，
用于统一处理 branch、jump、trap 和 \texttt{mret}；另一类是围绕 EX 的前递和停顿控制路径，
用于在保持顺序流水结构清晰的前提下收敛性能损失。

\subsection{数据相关处理}

为了在顺序流水中兼顾正确性和吞吐，项目实现了明确的数据相关处理策略：
\begin{itemize}
    \item EX/MEM、MEM/WB 两级前递，降低普通 ALU 相关导致的停顿；
    \item 针对 \texttt{load-use} 场景设置专门检测逻辑，当结果尚未可用且后继指令立即消费时暂停译码级；
    \item 译码级停顿条件采用 \texttt{stall\_decode = load\_use\_hazard} 策略，使停顿范围尽量收敛在必要场景内。
\end{itemize}

这一选择的直接效果是：对可前递场景尽量放行，对不可前递的 load-use 场景执行最小必要停顿，
在结构复杂度、稳定性和性能之间取得平衡。

\subsection{控制相关与 redirect 机制}

控制流处理主要在 EX 级完成决策：
\begin{itemize}
    \item branch、jump 和 \texttt{jalr} 在 EX 级形成 redirect；
    \item \texttt{trap}、\texttt{mret} 与普通跳转共享控制流重定向框架；
    \item 当 \texttt{ex\_fetch\_redirect\_valid} 有效时，IF 级更新 PC 并清空错误路径气泡。
\end{itemize}

这种组织方式能够用统一框架处理普通跳转、异常进入和异常返回，既保持控制路径简洁，
也便于验证 trap 与正常执行之间的切换关系。
赛题在 IF 级还提出了预取缓冲、BTB、分支预测器等更积极的前端机制。
当前正式实现采用 redirect/flush 与译码早重定向框架作为稳定基线，优先保证功能闭环、验证收敛和 FPGA 路径一致；
更完整的预取、BTB 与分支预测仍属于后续继续对齐赛题的增强方向，本文在此明确保留这一边界，而不对未实现能力作超额宣称。

\subsection{CSR、异常与中断处理}

本设计实现的控制状态寄存器与异常处理，重点覆盖比赛阶段确有需求的机器态功能：
\begin{itemize}
    \item CSR 指令：\texttt{csrrw/csrrs/csrrc} 及其立即数变体；
    \item 关键机器态 CSR：\texttt{mstatus}、\texttt{mie}、\texttt{mtvec}、\texttt{mscratch}、
    \texttt{mepc}、\texttt{mcause}、\texttt{mip}；
    \item 同步异常：非法指令、CSR 非法访问、load misaligned 等；
    \item 控制流相关：\texttt{ecall}、\texttt{ebreak}、\texttt{mret}；
    \item 机器定时器中断：通过 \texttt{timer} 外设触发并进入 trap 流程。
\end{itemize}

这一范围能够满足现阶段软件运行和验证需要，也有利于将实现重点集中在基础执行链路、异常响应和工程完整性上。

\subsection{SoC 存储与 MMIO 组织}

SoC 采用“ROM + RAM + MMIO”的最小闭环结构。其中同步 ROM 负责取指，RAM 负责数据访问，
MMIO 外设负责提供对外可观测接口和计时能力。关键地址组织如下：
\begin{table}[H]
    \centering
    \caption{SoC 关键地址映射}
    \small
    \begin{tabularx}{0.95\textwidth}{L{0.27\textwidth} L{0.23\textwidth} Y}
        \toprule
        资源 & 地址/基址 & 作用 \\
        \midrule
        ROM & \texttt{0x0000\_0000} & 程序存储基址 \\
        RAM & \texttt{0x0000\_4000} & 数据存储基址 \\
        UART TX & \texttt{0x1000\_0000} & 串口发送寄存器 \\
        DONE & \texttt{0x1000\_0004} & 程序结束标志寄存器 \\
        TIMER VALUE LO & \texttt{0x1000\_0008} & 计时值低 32 位 \\
        TIMER VALUE HI & \texttt{0x1000\_000C} & 计时值高 32 位 \\
        TIMER CMP LO & \texttt{0x1000\_0010} & 比较寄存器低 32 位 \\
        TIMER CMP HI & \texttt{0x1000\_0014} & 比较寄存器高 32 位 \\
        TIMER CTRL & \texttt{0x1000\_0018} & 定时器控制与中断使能 \\
        \bottomrule
    \end{tabularx}
\end{table}

这种组织方式具有三方面优势：一是能够直接支撑仿真、CoreMark 和 \texttt{riscv-tests}；
二是向 FPGA 原型迁移成本较低；三是 UART、timer 和 DONE 为后续 bring-up 提供了最小但有效的观测接口。

\subsection{参数化与扩展验证支撑}

虽然冻结提交口径以 \texttt{RV32I + Zicsr} 为主，但当前工程并没有把结构边界固化为单一 XLEN。
CPU 顶层、测试脚本和回归清单已经具备参数化与多清单组织能力，使项目能够在保持提交主口径稳定的同时，
完成 \texttt{RV32/RV64} 双 XLEN baseline 与扩展 \texttt{full-ui} 验证。这样的工程组织有两项直接收益：
\begin{itemize}
    \item 一方面，它说明当前设计不是“只能跑单一路径”的一次性实验实现，而是具备继续扩展和继续验证的基础；
    \item 另一方面，它使冻结口径、扩展验证口径和后续优化工作能够被清晰地区分，避免结果口径漂移。
\end{itemize}

\subsection{架构取舍}

本设计优先采用边界清晰、易于验证和便于复现的结构，重点保留已经通过验证矩阵的机制与参数，
暂不引入 cache、多发射、乱序等更高风险的复杂结构。这样的取舍使得文中的关键结论能够被实验结果直接支撑，
同时也为后续性能优化和实体板验证保留了清晰起点。
